\documentclass[11pt]{letter}

\usepackage[margin=0.75in]{geometry}
\usepackage{setspace}
\usepackage{hyperref}
\usepackage{xcolor}
\usepackage{enumitem}
\usepackage{parskip}

\hypersetup{
    colorlinks=true,
    urlcolor=black,
    linkcolor=black
}

\setlength{\parindent}{0pt}
\setstretch{1.05}
\date{}
\begin{document}

\begin{letter}{}

\opening{Dear Hiring Committee,}

I saw the job post on \textbf{Amazon's careers site} and thought I'd be a good fit for it.

I am a Master's student in Computer Science at the University of Cincinnati, graduating December 2026, looking for a role as a \textbf{Software Engineer, AI/ML} on the \textbf{Annapurna Labs} Distributed Training team at \textbf{Amazon}. With hands on experience automating large-scale AWS infrastructure and building PyTorch-based performance evaluation pipelines, I am excited about the opportunity to contribute to \textbf{AWS Neuron} and help push training throughput on Trainium.

\textbf{Here are reasons that make me a great fit for the role:}

\begin{itemize}[leftmargin=*, itemsep=0.45em]
    \item \textbf{Distributed systems and performance profiling:} As a Graduate Researcher, I built a PyTorch-based evaluation pipeline that profiled a two-stage LLM system end to end to isolate a per-step inference bottleneck, achieving a 25.7x latency reduction (39.4s vs 1012.1s avg per run) across 53 runs on 20+ Android apps while preserving task-completion fidelity.

    \item \textbf{Large-scale AWS infrastructure experience:} At Siemens Industry Software, I automated Teamcenter X infrastructure workflows with Python (Boto3), cutting deployment time by 40\% across AWS environments, and executed 44 production environment upgrades in 4 months, diagnosing live infrastructure issues across Linux, EC2, and databases.

    \item \textbf{System design across the stack:} I migrated 40+ production environments from process-based deployments to Kubernetes using ArgoCD, Helm, and GitLab CI/CD, and designed an Ansible-based Blue-Green deployment framework that reduced production upgrade downtime by 50\%, giving me practical experience coordinating changes across interdependent infrastructure layers.

    \item \textbf{Engineering quality and collaboration:} As Quality Coach and SonarQube Specialist for the Agile Release Train, I increased automated code coverage from 35\% to 78\% and mentored 5 software engineers, reflecting the cross-team collaboration this role requires when tracing performance gaps across frameworks, compiler, and runtime teams.
\end{itemize}

I am available for full-time work after my graduation in December 2026. As a Master's student at the University of Cincinnati and someone who values professionalism and growth, I am enthusiastic about contributing to Neuron's mission of accelerating deep learning and GenAI workloads on Trainium.

Please find my CV attached below. I look forward to hearing from you. I would be excited to work at the intersection of ML and distributed systems, helping shape the direction of AI acceleration technology.

{\setlength{\parskip}{0pt}
Regards,\\
Tanmay Bhuskute,\\
bhuskutv@mail.uc.edu $\mid$ tanmay.v.bhuskute@gmail.com $\mid$ +1 (732) 858-3027,\\
3424 Brookline Ave, Flat 19, Cincinnati, OH 45220.
}

\end{letter}

\end{document}
